%% -*- coding: utf-8 -*-
%% ROLE: engine-vendored — DO NOT edit here. `scitex-writer update-project`
%% overwrites this file on every re-vendor; fix it upstream in the
%% scitex-writer package instead (local edits are lost, and update-project
%% may set it read-only in the consumer workspace after vendoring).
%% File: 00_shared/latex_styles/formatting.tex

%% Keep the shared title available for delayed elsarticle page headers.
%% ROLE: engine-vendored — preserve consumer-owned 00_shared/title.tex.
%% elsarticle retains the shared title for page headers after frontmatter ends.
%% A \newcommand defined by the title input inside that environment is local;
%% export its existing value before the group closes, without redefining text.
\AddToHook{env/frontmatter/end}{%
  \ifdefined\scitexmanuscripttitle
    \global\let\scitexmanuscripttitle\scitexmanuscripttitle
  \fi
}


%% --- Journal name header (show journal name as-is, without "Preprint submitted to" prefix) ---
\makeatletter
\def\ps@pprintTitle{%
  \def\@oddhead{\reset@font\footnotesize\itshape\@journal\hfill}%
  \let\@evenhead\@oddhead
  \def\@oddfoot{\reset@font\hfil\thepage\hfil}%
  \let\@evenfoot\@oddfoot}
\makeatother

%% --- Image width ---
\newlength{\imagewidth}
\newlength{\imagescale}

%% --- Line numbers ---
\linespread{1.2}
\linenumbers

%% --- Float placement (reduce stranded whitespace from large/wide floats) ---
%% LaTeX defaults push a big table/figure onto a near-empty page, leaving the
%% rest of that page blank. These let floats fill pages more fully before a
%% float-only page is started, and let a float-only page hold a single large
%% float without demanding it be nearly full.
\renewcommand{\topfraction}{0.9}        % a top float may cover up to 90% of a page
\renewcommand{\bottomfraction}{0.75}
\renewcommand{\textfraction}{0.1}       % keep >=10% text before forcing a float page
\renewcommand{\floatpagefraction}{0.7}  % a float-only page must be >=70% full
\setcounter{topnumber}{3}
\setcounter{bottomnumber}{2}
\setcounter{totalnumber}{5}

%% --- Colors ---
\definecolor{GreenBG}{rgb}{0,1,0}
\definecolor{RedBG}{rgb}{1,0,0}

%% --- SciTeX theme colors (source: figrecipe/colors/_PARAMS.py) ---
\definecolor{SciTeXBlue}{HTML}{0080C0}
\definecolor{SciTeXRed}{HTML}{FF4632}
\definecolor{SciTeXPink}{HTML}{FF96C8}
\definecolor{SciTeXGreen}{HTML}{14B414}
\definecolor{SciTeXYellow}{HTML}{E6A014}
\definecolor{SciTeXGray}{HTML}{808080}
\definecolor{SciTeXPurple}{HTML}{C832FF}
\definecolor{SciTeXLightBlue}{HTML}{14C8C8}
\definecolor{SciTeXBrown}{HTML}{800000}
\definecolor{SciTeXNavy}{HTML}{000064}
\definecolor{SciTeXOrange}{HTML}{E45E32}

%% --- Highlight boxes ---
\newtcbox{\greenhighlight}[1][]{on line,colframe=GreenBG,colback=GreenBG!50!white,boxrule=0pt,arc=0pt,boxsep=0pt,left=1pt,right=1pt,top=2pt,bottom=2pt,tcbox raise base}
\newtcbox{\redhighlight}[1][]{on line,colframe=RedBG,colback=RedBG!50!white,boxrule=0pt,arc=0pt,boxsep=0pt,left=1pt,right=1pt,top=2pt,bottom=2pt,tcbox raise base}

\newcommand{\REDSTARTS}{\color{red}}
\newcommand{\REDENDS}{\color{black}}
\newcommand{\GREENSTARTS}{\color{green}}
\newcommand{\GREENENDS}{\color{black}}

%% --- Word count ---
%% A word-count file is decorative. It may be MISSING (not yet generated, or the
%% path is CWD-relative and pdflatex ran from another dir). Guard the read: a
%% missing/unopened stream must NOT reach a bare \read, which would block on the
%% terminal and, under -interaction=nonstopmode, EMERGENCY-STOP the whole
%% compile ("<read N>") -> no PDF -> empty .aux -> bibtex finds no \citation ->
%% every \cite undefined (??/[?]). \ifeof after \openin is true when the file
%% could not be opened (or is empty); default to 0 + a loud \typeout warning and
%% keep compiling.
\newread\wordcount
\newcommand\readwordcount[1]{%
\def\thewordcount{0}%
\openin\wordcount=#1\relax
% Guard: a missing OR empty count file must NEVER reach \num as an empty
% argument -- siunitx then errors "Invalid number", which is FATAL and fails
% every latexmk pass (so citations render as [?]). \ifeof is true for both a
% missing file and a 0-byte file; the post-read \ifx catches a blank line.
% Default to 0 in all those cases. (neurovista incident 2026-07-04.)
\ifeof\wordcount
  \typeout{WARNING: scitex-writer: word-count file '#1' not found or empty; using 0.}%
\else
  \read\wordcount to \thewordcount
\fi
\closein\wordcount
\ifx\thewordcount\empty\def\thewordcount{0}\fi
\begingroup\sisetup{round-mode=none,group-separator={,},group-minimum-digits=4}\num{\thewordcount}\endgroup%
}

%% --- Text highlighting ---
\usepackage{soul}
\sethlcolor{yellow}

%% --- Reference handling ---
\usepackage{refcount}

\let\oldref\ref
\newcommand{\hlref}[1]{%
  \ifnum\getrefnumber{#1}=0
    \colorbox{yellow}{\ref*{#1}}%  % Use colorbox for references (no line break needed)
  \else
    \ref{#1}%
  \fi
}

% To add an 'S' prefixes to a reference
\newcommand*\sref[1]{S\hlref{#1}}
\newcommand*\sfref[1]{Supplementary Figure S\hlref{#1}}
\newcommand*\stref[1]{Supplementary Table S\hlref{#1}}
\newcommand*\smref[1]{Supplementary Materials S\hlref{#1}}

%% --- Placeholder for draft content ---
\newcommand{\placeholder}[1]{\colorbox{yellow}{\textcolor{red}{\textbf{[#1]}}}}
\newcommand{\ph}[1]{\placeholder{#1}}

%% --- Controlled inline figure placement ---
%% By default figures collect in the end "Figures" section. To place a figure
%% in the main text at a controlled spot, drop \scitexfig{<number>} (e.g.
%% \scitexfig{01}) where you want it: it renders that figure's float here and
%% flags it so it is NOT also repeated in the end section. Unplaced figures
%% still collect at the end (default behaviour preserved). \scitexfigdir points
%% at the assembler-generated per-figure files; override it for non-manuscript
%% doctypes. The flag csname uses '@' but is built via \csname (catcode
%% irrelevant -> flatten-safe; no \makeatletter needed).
\providecommand{\scitexfigdir}{./01_manuscript/contents/figures/compiled/_placeable}
\newcommand{\scitexfig}[1]{%
  \expandafter\gdef\csname scitexfigplaced@#1\endcsname{}%
  \input{\scitexfigdir/#1}%
}

%% --- Controlled inline table placement (mirror of \scitexfig) ---
%% Drop \scitextab{<number>} (e.g. \scitextab{01}) in the body to render that
%% table here; it is then skipped in the end "Tables" section. Unplaced tables
%% still collect at the end. Override \scitextabdir for non-manuscript doctypes.
\providecommand{\scitextabdir}{./01_manuscript/contents/tables/compiled/_placeable}
\newcommand{\scitextab}[1]{%
  \expandafter\gdef\csname scitextabplaced@#1\endcsname{}%
  \input{\scitextabdir/#1}%
}

%% --- Inline citation-verification render (clew) ---
%% \clewverified{<cite text>}{<status>}{<source url>}: per-claim trust signal
%% sharing clew's verify_citations result with the page-1 banner. A "verified"
%% status renders the text underlined + hyperlinked to its source; any other
%% status (unverified/stub/unknown) renders bare -- "unverified = no
%% decoration". Provided (not new) so it is a safe no-op default before the
%% clew wiring lands and callers may already define their own. No '@' catcodes,
%% so it survives flattening without \makeatletter.
\providecommand{\clewverifiedstatus}{verified}
\providecommand{\clewverified}[3]{%
  \def\clewgivenstatus{#2}%
  \ifx\clewgivenstatus\clewverifiedstatus
    \href{#3}{\ul{#1}}%
  \else
    #1%
  \fi
}

%%%% EOF